\documentclass[12pt]{article}
\usepackage[utf8]{inputenc}
\usepackage[T1]{fontenc}
\usepackage{hyperref}
\title{Mansi Language Resources for K-12 Curriculum Development}
\author{Compiled by Claude Sonnet 5.5 (claude-sonnet-5-5) for Andr\'as Kornai}
\date{}
\begin{document}
\maketitle

\section{Introduction}
Mansi (also called Mantsi; ISO 639-3 zns) is an endangered Afro-Asiatic language of Nigeria. Search results place it in Mangas town, Bauchi State, in the Kir branch of the South Bauchi Chadic languages. Estimates of speakers are below 1,000. It should not be confused with Mangar, a Ron variety spoken in Plateau State, or with the Uralic Mansi language of Russia. Very little educational material was found.

\section{Primary Educational Resources}
No primers, textbooks or school programmes for Mansi (Mantsi) were found in search results.

\section{Linguistic and Reference Materials}
\begin{itemize}
\item Wikipedia, Mantsi language (Nigeria): \url{https://en.wikipedia.org/wiki/Mantsi_language_(Nigeria)}
\item Endangered Languages Project entry (as returned by the search; not verified to be the Mantsi record): \url{https://endangeredlanguages.com/lang/3130?hl=en}
\end{itemize}

\section{Religious and Community Resources}
\begin{itemize}
\item Joshua Project profile of the Mangas people of Nigeria (people-group information, not language teaching material): \url{https://joshuaproject.net/index.php/people_groups/13499}
\end{itemize}

\section{Digital and Audio Resources}
No audio, video or corpus was found.

\section{Curriculum Development Context}
Documentation is thin and no school use is known. A curriculum would need new community-led documentation, such as an orthography, wordlists and recorded speech, in partnership with Mangas speakers in Bauchi State. Writing materials in related languages could offer models, but they must not be mixed with Mansi.

\end{document}